\section{Descrição do Processo}\label{sec:processo}

\subsection{Configuração Geral da UTAA}

O esquema de Tratamento de Águas Ácidas considerada neste trabalho é o esquema composto por duas torres de destilação operando em série.

Nesse esquema de tratamento, a água ácida proveniente de diversas unidades da refinaria (hidrotratamento, FCC, coqueamento, etc.) é coletada e alimentada na primeira torre, a torre esgotadora de H\textsubscript{2}S. Esta coluna opera sob pressão moderada e utiliza vapor de esgotamento e um refervedor para promover a evaporação preferencial do H\textsubscript{2}S, que é mais volátil que a NH\textsubscript{3} nas condições operacionais típicas. O gás ácido de topo, rico em H\textsubscript{2}S e contendo alguma NH\textsubscript{3}, é enviado para a URE. A água de fundo, ainda contendo NH\textsubscript{3} e traços de H\textsubscript{2}S, é bombeada para a segunda torre, onde a NH\textsubscript{3} é removida.


\subsection{Torre Esgotadora de H2S}

A torre esgotadora de H\textsubscript{2}S é o foco principal deste trabalho. Trata-se de uma coluna de destilação com as seguintes características típicas:

\begin{itemize}[leftmargin=*]
  \item Número de estágios: 30 
  \item Eficiência de Murphree (em toda a coluna): 70\%
  \item Pressão de operação: [valor] bar
  \item Alimentação: água ácida a [temperatura] e [pressão]
  \item Vapor de despojamento: injetado no fundo da coluna
  \item Refervedor: tipo [kettle/termossifão], carga térmica [valor] kW
  \item Produto de topo: gás ácido para URE
  \item Produto de fundo: água parcialmente tratada para segunda torre
\end{itemize}

\subsection{Desafios Operacionais}

A operação da torre esgotadora apresenta os seguintes desafios:

\begin{enumerate}
  \item \textbf{Maximização da recuperação de H\textsubscript{2}S}: É desejável maximizar a concentração de H\textsubscript{2}S no gás ácido para otimizar a operação da URE.
  \item \textbf{Minimização de NH\textsubscript{3} no gás ácido}: A presença de NH\textsubscript{3} no gás ácido leva à formação de sais de amônio (principalmente NH\textsubscript{4}HS), que causam corrosão, entupimentos e redução de eficiência na URE.
  \item \textbf{Prevenção de fervura da coluna}: Em condições extremas de carga térmica ou composição, pode ocorrer o fenômeno de "fervura", caracterizado por um aumento abrupto da taxa de vaporização que leva a arraste excessivo de líquido e NH\textsubscript{3} para o topo.
  \item \textbf{Eficiência energética}: A carga térmica do refervedor representa um custo operacional significativo, sendo desejável minimizá-la mantendo os objetivos de separação.
\end{enumerate}


% TODO: Detalhar parâmetros do GPSWAT utilizados e reações consideradas

\subsubsection{Especificações do Processo}

As principais especificações do modelo estático incluem:

\begin{table}[H]
  \centering
  \caption{Especificações do modelo estático da torre esgotadora de H\textsubscript{2}S.}
  \label{tab:spec_estatico}
  \begin{tabular}{lcc}
    \toprule
    \textbf{Parâmetro} & \textbf{Valor} & \textbf{Unidade} \\
    \midrule
    Número de estágios & [valor] & - \\
    Estágio de alimentação & [valor] & - \\
    Pressão de topo & [valor] & bar \\
    Pressão de fundo & [valor] & bar \\
    Vazão de alimentação & [valor] & kg/h \\
    Temperatura de alimentação & [valor] & \si{\celsius} \\
    Vazão de vapor de despojamento & [valor] & kg/h \\
    Razão de refluxo & [valor] & - \\
    Carga térmica do refervedor & [valor] & kW \\
    \bottomrule
  \end{tabular}
\end{table}

% TODO: Preencher valores após definição do caso base

\subsubsection{Composição da Alimentação}

A composição típica da água ácida de alimentação é apresentada na Tabela~\ref{tab:composicao_feed}.

\begin{table}[H]
  \centering
  \caption{Composição típica da alimentação de água ácida.}
  \label{tab:composicao_feed}
  \begin{tabular}{lcc}
    \toprule
    \textbf{Componente} & \textbf{Concentração} & \textbf{Unidade} \\
    \midrule
    H\textsubscript{2}O & [valor] & wt\% \\
    H\textsubscript{2}S & [valor] & ppm (w) \\
    NH\textsubscript{3} & [valor] & ppm (w) \\
    CO\textsubscript{2} & [valor] & ppm (w) \\
    Outros & [valor] & ppm (w) \\
    \bottomrule
  \end{tabular}
\end{table}

% TODO: Definir composição de referência e variações para estudo paramétrico



\subsection{Modelagem Dinâmica no Aspen Plus Dynamics V14}\label{sec:modelo_dinamico}

\subsubsection{Preparação do Modelo Estático}

Antes da exportação para o ambiente dinâmico, o modelo estático foi preparado conforme os seguintes passos:

\begin{enumerate}
  \item Verificação de convergência robusta do modelo estático
  \item Dimensionamento de equipamentos (diâmetro de coluna, volumes de pratos, holdups)
  \item Especificação de controladores básicos de inventário (nível, pressão)
  \item Definição de válvulas de controle e suas características
  \item Configuração de instrumentação (transmissores, sensores)
\end{enumerate}

\subsubsection{Exportação e Configuração em Aspen Dynamics}

O modelo foi exportado do Aspen Plus para o Aspen Plus Dynamics utilizando a funcionalidade de exportação de modelos em regime permanente (\emph{Pressure-Driven} mode). As principais configurações incluem:

\begin{itemize}[leftmargin=*]
  \item Modo de simulação: modo orientado por pressão (\emph{pressure-driven})
  \item Passo de integração: [valor] segundos
  \item Método de integração: [especificar]
  \item Configuração de controladores PID para inventários
  \item Especificação de dinâmicas de válvulas
\end{itemize}

\subsubsection{Condições Iniciais e Convergência Dinâmica}

A inicialização do modelo dinâmico utilizou o estado estacionário convergido do modelo estático. As seguintes verificações foram realizadas:

\begin{enumerate}
  \item Balanços de massa e energia em cada estágio
  \item Estabilidade de holdups e taxas de fluxo
  \item Resposta de controladores de inventário
  \item Tempo de estabelecimento após pequenas perturbações
\end{enumerate}

\subsection{Cenários de Simulação e Análise}\label{sec:cenarios}

Para investigar o comportamento do sistema em torno do ponto crítico de operação (limite de fervura), foram definidos os seguintes cenários:

\begin{enumerate}
  \item \textbf{Caso Base}: Condições nominais de operação com composição de alimentação típica
  \item \textbf{Variação de H\textsubscript{2}S}: Aumento/diminuição da concentração de H\textsubscript{2}S na alimentação em [valor]\%
  \item \textbf{Variação de NH\textsubscript{3}}: Aumento/diminuição da concentração de NH\textsubscript{3} na alimentação em [valor]\%
  \item \textbf{Variação de carga térmica}: Aumento/diminuição da carga do refervedor em [valor]\%
  \item \textbf{Variação de vazão}: Aumento/diminuição da vazão de alimentação em [valor]\%
  \item \textbf{Aproximação do limite de fervura}: Condições operacionais progressivamente mais severas até atingir o limite
\end{enumerate}

\subsection{Métricas de Desempenho}

Para avaliar o desempenho operacional da torre em cada cenário, as seguintes métricas serão calculadas:

\begin{itemize}[leftmargin=*]
  \item \textbf{Recuperação de H\textsubscript{2}S}: fração mássica de H\textsubscript{2}S da alimentação que sai no gás ácido de topo
  \item \textbf{Concentração de NH\textsubscript{3} no gás ácido}: ppm (mol) de NH\textsubscript{3} no topo
  \item \textbf{Consumo energético}: carga térmica do refervedor por unidade de massa de água ácida tratada
  \item \textbf{Margem até fervura}: distância operacional até condições críticas
  \item \textbf{Qualidade da água tratada}: concentração residual de H\textsubscript{2}S e NH\textsubscript{3} no fundo
\end{itemize}

\section{Resultados e Discussão}\label{sec:resultados}

% TODO: Apresentar resultados das simulações estáticas e dinâmicas

\subsection{Modelo Estático: Caso Base}

[Apresentar resultados do caso base em regime permanente, perfis de composição, temperatura e pressão ao longo da coluna]

% \begin{figure}[H]
%   \centering
%   % \includegraphics[width=0.8\linewidth]{perfil_temperatura.png}
%   \caption{Perfil de temperatura ao longo da torre esgotadora de H\textsubscript{2}S — caso base.}
%   \label{fig:perfil_temp}
% \end{figure}

% \begin{figure}[H]
%   \centering
%   % \includegraphics[width=0.8\linewidth]{perfil_composicao.png}
%   \caption{Perfis de composição de H\textsubscript{2}S e NH\textsubscript{3} ao longo da coluna — caso base.}
%   \label{fig:perfil_comp}
% \end{figure}

\subsection{Análise de Sensibilidade Paramétrica}

[Apresentar resultados da análise de sensibilidade para variações de composição, carga térmica e vazão]

\subsubsection{Efeito da Concentração de H2S na Alimentação}

[Discutir como variações de H\textsubscript{2}S afetam recuperação, consumo energético e operação]

\subsubsection{Efeito da Concentração de NH3 na Alimentação}

[Discutir como variações de NH\textsubscript{3} afetam o arraste de NH\textsubscript{3} para o topo e o risco de fervura]

\subsubsection{Efeito da Carga Térmica do Refervedor}

[Analisar trade-offs entre carga térmica, recuperação de H\textsubscript{2}S e arraste de NH\textsubscript{3}]

\subsection{Modelo Dinâmico: Resposta Transitória}

[Apresentar resultados de simulações dinâmicas para perturbações típicas]

% \begin{figure}[H]
%   \centering
%   % \includegraphics[width=0.8\linewidth]{resposta_dinamica.png}
%   \caption{Resposta dinâmica a perturbação em degrau na composição de alimentação.}
%   \label{fig:resposta_dinamica}
% \end{figure}

\subsection{Identificação do Limite de Fervura}

[Apresentar caracterização do ponto crítico de operação, sinais precursores de fervura, margem operacional segura]

\subsection{Discussão dos Resultados}

[Sintetizar principais descobertas, comparar com dados de literatura, discutir implicações práticas para operação de refinaria]

\FloatBarrier

\section{Conclusões}\label{sec:conclusoes}

% TODO: Sintetizar principais conclusões do trabalho

Este trabalho apresentou o desenvolvimento e análise de modelos de simulação estáticos e dinâmicos de uma Unidade de Tratamento de Águas Ácidas (UTAA), com foco particular na torre esgotadora de H\textsubscript{2}S. Os modelos foram implementados em Aspen Plus V14 e Aspen Plus Dynamics V14, utilizando o pacote termodinâmico GPSWAT para representação adequada dos equilíbrios físico-químicos envolvendo eletrólitos em solução aquosa.

[Discutir principais conclusões sobre:]
\begin{itemize}[leftmargin=*]
  \item Capacidade dos modelos de reproduzir o comportamento esperado do processo
  \item Trade-offs identificados entre recuperação de H\textsubscript{2}S, minimização de NH\textsubscript{3} e eficiência energética
  \item Caracterização do ponto operacional crítico (limite de fervura)
  \item Sensibilidade do sistema a variações de composição e condições operacionais
  \item Implicações para estratégias de operação e controle
\end{itemize}

Os modelos desenvolvidos fornecem uma base sólida para otimização operacional da UTAA e para o desenvolvimento de estratégias de controle avançado, conforme discutido na seção de trabalhos futuros.

\section{Trabalhos Futuros}\label{sec:futuros}

Este trabalho estabelece as fundações para desenvolvimentos subsequentes em controle avançado da UTAA. As seguintes direções são propostas:

\subsection{Geração de Dados para Treinamento de Redes Neurais}

O modelo dinâmico desenvolvido em Aspen Plus Dynamics será utilizado para gerar extensos conjuntos de dados de treinamento mediante simulação de trajetórias operacionais diversificadas. Especificamente:

\begin{itemize}[leftmargin=*]
  \item Simulação de múltiplos cenários com variações nas composições de alimentação, vazões e condições operacionais
  \item Coleta de séries temporais de variáveis de processo (temperaturas, composições, pressões, fluxos)
  \item Identificação de relações entrada-saída para mapeamento via redes neurais
  \item Geração de dados próximos ao limite de fervura para capturar não-linearidades críticas
\end{itemize}

\subsection{Desenvolvimento de Modelos de Redes Neurais}

Com base nos dados gerados, redes neurais artificiais (ANNs) serão treinadas para emular o comportamento dinâmico da torre esgotadora. Arquiteturas candidatas incluem:

\begin{itemize}[leftmargin=*]
  \item Redes neurais feedforward (FNNs) para mapeamento estático entrada-saída
  \item Redes neurais recorrentes (RNNs, LSTMs, GRUs) para captura de dinâmicas temporais
  \item Modelos híbridos combinando conhecimento fenomenológico com componentes baseados em dados
\end{itemize}

\subsection{Controle Preditivo Não Linear Baseado em Redes Neurais (NN NMPC)}

O objetivo final é desenvolver um controlador preditivo não linear (NMPC) utilizando o modelo de rede neural como preditor interno. A estrutura proposta envolve:

\begin{enumerate}
  \item \textbf{Modelo de predição}: Rede neural treinada substituindo o modelo fenomenológico rigoroso no otimizador do NMPC, reduzindo drasticamente o custo computacional
  \item \textbf{Função objetivo}: Minimização de um funcional custo considerando rastreamento de referências (recuperação de H\textsubscript{2}S), penalização de NH\textsubscript{3} no topo, esforço de controle e eficiência energética
  \item \textbf{Restrições}: Limites operacionais em variáveis manipuladas (carga térmica, vazão de vapor) e variáveis de processo (temperatura, pressão, concentrações)
  \item \textbf{Estimação de estados}: Filtros estendidos de Kalman (EKF) ou estimadores por horizonte móvel (MHE) para estados não medidos
  \item \textbf{Operação offset-free}: Incorporação de estados integradores para rejeição de distúrbios e desalinhamentos modelo-planta
\end{enumerate}

\subsection{Validação e Implementação}

O controlador NN NMPC será validado inicialmente em simulações fechadas (closed-loop) utilizando o modelo rigoroso do Aspen Plus Dynamics como "planta virtual". Serão avaliados:

\begin{itemize}[leftmargin=*]
  \item Desempenho de rastreamento de setpoints
  \item Rejeição de distúrbios em composição e vazão de alimentação
  \item Capacidade de operação próxima ao limite de fervura sem violação de restrições
  \item Robustez a incertezas de modelo e ruídos de medição
  \item Comparação com estratégias de controle convencionais (PID, MPC linear)
\end{itemize}

Caso os resultados sejam promissores, estudos de viabilidade para implementação em planta piloto ou industrial poderão ser conduzidos.
